\documentclass[uplatex,dvipdfmx,a4paper,11pt]{jsarticle}
\input{../../../templates/kaisetsu.tex}
\kaisetsuprint{量子力学1　第2章「シュレーディンガー方程式と対応原理」}{2026-10-29}
%% kaisetsu-moto	problems\量子力学1\shou02\量子力学1-2.tex	kiroku\kaisetsu\量子力学1\量子力学1-2_2026-10-29_komento.tex
\begin{document}
\matome
\begin{saitenichiran}
  \hyougyou{\hyperlink{bt-m-016-stationary}{5}}{L3・導出}{◎△◎}{B}{(2)方針}
\end{saitenichiran}
\begin{daiji}
  \item 位相速度は1つの平面波の位相が一定の点の速さ $\omega/k$ で、自由粒子では $\hbar k/(2m)$、つまり粒子の速さ $\hbar k/m$ の半分である。粒子の速さに当たるのは群速度 $\dv{\omega}{k}=\hbar k/m$ のほうで、両者を混同しない（問題5 (2)）。
  \item 波の山の速さは「位相（差）が一定」の式を書き、$x$ を $t$ で表して出す。2つの平面波の重ね合わせでは差の比 $(\omega_{2}-\omega_{1})/(k_{2}-k_{1})$ で、微分ではない（問題5 (2)）。
\end{daiji}
\begin{maemo}
  \item 位相一定の根拠の書き落とし：量子力学1-2 問題1（平面波の位相速度，2026-10-14）の (1) でも、位相が一定になる点を追うという根拠を書いていなかった。速さを問われたら、まず「位相 $=$ 一定」の式を1行書いてから $x$ を $t$ で表す。
\end{maemo}

\mondaihen
\level{3}
\mondaibangou{5}
\begin{mondai}{016-stationary}{導出}{時間を含むシュレーディンガー方程式と平面波の分散関係，重ね合わせた波の山の速さ（群速度）と定常状態の重ね合わせの時間変化}
\begin{teisuu}
  プランク定数 & $h$ & \qty{6.626e-34}{\joule\second}\\
  換算プランク定数 & $\hbar$ & \qty{1.0546e-34}{\joule\second}\\
  電子の質量 & $\melec$ & \qty{9.109e-31}{\kilogram}\\
  電気素量 & $e$ & \qty{1.602e-19}{\coulomb}\\
  真空中の光速 & $c$ & \qty{2.998e8}{\metre\per\second}\\
\end{teisuu}
この問題では，要点整理の時間を含むシュレーディンガー方程式を用いる。1次元で考える。
  \begin{shoumon}
    \item 自由粒子（$V=0$）で $\psi_{k}(x,t)=\exp[\ii(kx-\omega t)]$ が時間を含むシュレーディンガー方程式の解となるための $\omega$ と $k$ の関係を求めよ。また，$E=\hbar\omega$，$p=\hbar k$ を用いてその意味を述べよ。
    \item $0<k_{1}<k_{2}$ とし，$\omega_{1}$，$\omega_{2}$ は (1) の関係を満たすとして，$\psi=\psi_{k_{1}}+\psi_{k_{2}}$ を考える。$\abs{\psi}^{2}$ を $x$ と $t$ の実関数で表し，$\abs{\psi}^{2}$ の山が動く速さ $v_{\mathrm{g}}$ を $k_{1}$，$k_{2}$，$m$，$\hbar$ で表せ。$v_{\mathrm{g}}$ を，2つの平面波に対応する粒子の速さ $\hbar k_{1}/m$，$\hbar k_{2}/m$ と比べよ。さらに，$k_{1}$ と $k_{2}$ が近いとき，$v_{\mathrm{g}}$ は平面波1つの位相速度の何倍かを答えよ。
    \item 幅 $L$ の無限に深い井戸（要点整理）の $\phi_{n}(x)=\sin(n\pi x/L)$ を用いて，$\psi_{n}(x,t)=\phi_{n}(x)\ee^{-\ii E_{n}t/\hbar}$ とおく。$\abs{\psi_{n}}^{2}$ が時間によらないことを示せ。また，$\psi=\psi_{1}+\psi_{2}$ について $\abs{\psi}^{2}$ を求め，$\abs{\psi}^{2}$ が時間変化する角振動数を求めよ。この角振動数を $2\pi$ で割った振動数と，粒子が $n=2$ から $n=1$ の状態に移るときに放出される光子の振動数との関係を述べよ。
  \end{shoumon}

\begin{kotaeran}
  \machigai{2}{△}{方針} $v_{\mathrm{g}}=\hbar(k_{1}+k_{2})/(2m)$ と「$\hbar k_{1}/m$ と $\hbar k_{2}/m$ の平均」は正しい。最後の問いに「$k_{1}=k_{2}$ のとき $v_{\mathrm{g}}=\hbar k/m$，1倍」と答えているが、平面波1つの位相速度は (1) の $\omega=\hbar k^{2}/(2m)$ から $\omega/k=\hbar k/(2m)$ で、$\hbar k/m$ ではない（$\hbar k/m$ は粒子の速さ）。正しくは $v_{\mathrm{g}}\approx\hbar k/m=2\times\omega/k$ で約2倍。

  ほかに2点：$\abs{\psi}^{2}$ を $2+\exp[\ii((k_{2}-k_{1})x-(\omega_{2}-\omega_{1})t)]+\exp[\ii((k_{1}-k_{2})x-(\omega_{1}-\omega_{2})t)]$ のままにしているが、問いは実関数で表すことなので $2+2\cos[(k_{2}-k_{1})x-(\omega_{2}-\omega_{1})t]$ までまとめる。また $v_{\mathrm{g}}=\dd(\omega_{2}-\omega_{1})/\dd(k_{2}-k_{1})$ と微分の形で書いているが、山の速さは $\cos$ の中身が一定の点の動き $x=[(\omega_{2}-\omega_{1})/(k_{2}-k_{1})]\,t+\text{定数}$ から出る差の比である（計算は実際には割り算になっていて、結果は合っている）。
\end{kotaeran}
\begin{rikai}
  \item 位相速度：1つの平面波 $\ee^{\ii(kx-\omega t)}$ で $kx-\omega t=$ 一定の点は $x=(\omega/k)t+$ 定数 と動くので、速さは $\omega/k$。自由粒子では $\omega/k=\hbar k/(2m)=p/(2m)$ で、粒子の速さ $p/m$ の半分になる。
  \item 群速度：2つの波の交差項の位相 $(k_{2}-k_{1})x-(\omega_{2}-\omega_{1})t$ が一定の点の速さ $(\omega_{2}-\omega_{1})/(k_{2}-k_{1})$。$k_{2}\to k_{1}$ の極限が $\dv{\omega}{k}=\hbar k/m$ で、これが粒子の速さに一致する。
  \item 「実関数で表せ」と問われたら、複素共役どうしの和を $\ee^{\ii\theta}+\ee^{-\ii\theta}=2\cos\theta$ でまとめる。
\end{rikai}
\end{mondai}

\kaitouhen
\level{3}
\kaisetsubangou{5}
\begin{kaisetsu}{016-stationary}
\houshin 平面波を時間を含むシュレーディンガー方程式に代入して $\omega$ と $k$ の関係を得る。重ね合わせの $\abs{\psi}^{2}$ は交差項に位相の差が残るので，その位相の差が一定の点の動きから山の速さを求める。定常状態の重ね合わせでは，交差項にエネルギーの差に比例する角振動数が残る。
\keisan
\shoumonkai{1} $\psi_{k}$ を代入する。
\begin{align*}
  \text{左辺}&=\ii\hbar\pdv{\psi_{k}}{t}=\ii\hbar(-\ii\omega)\psi_{k}=\hbar\omega\,\psi_{k},\\
  \text{右辺}&=-\frac{\hbar^{2}}{2m}\pdv[2]{\psi_{k}}{x}=\frac{\hbar^{2}k^{2}}{2m}\psi_{k}.
\end{align*}
$\psi_{k}\neq0$ なので $\hbar\omega=\hbar^{2}k^{2}/(2m)$，すなわち $\omega=\hbar k^{2}/(2m)$。$E=\hbar\omega$，$p=\hbar k$ を使うと $E=p^{2}/(2m)$ で，自由粒子のエネルギーと運動量の関係を表す。

\shoumonkai{2} $\theta_{j}=k_{j}x-\omega_{j}t$ とおくと $\psi=\ee^{\ii\theta_{1}}+\ee^{\ii\theta_{2}}$ で，
\[ \abs{\psi}^{2}=1+1+\ee^{\ii(\theta_{2}-\theta_{1})}+\ee^{-\ii(\theta_{2}-\theta_{1})}=2+2\cos\left[(k_{2}-k_{1})x-(\omega_{2}-\omega_{1})t\right] \]
山は位相の差 $\theta_{2}-\theta_{1}$ が $2\pi$ の整数倍の点で，その点は $x=(\omega_{2}-\omega_{1})t/(k_{2}-k_{1})+\text{定数}$ と動く。(1) の関係を使うと
\[ v_{\mathrm{g}}=\frac{\omega_{2}-\omega_{1}}{k_{2}-k_{1}}=\frac{\hbar(k_{2}^{2}-k_{1}^{2})}{2m(k_{2}-k_{1})}=\frac{\hbar(k_{1}+k_{2})}{2m} \]
これは2つの粒子の速さ $\hbar k_{1}/m$ と $\hbar k_{2}/m$ の平均である。$k_{1}\approx k_{2}\approx k$ のとき $v_{\mathrm{g}}\approx\hbar k/m$ で，これは要点整理の群速度 $\dv{\omega}{k}=\hbar k/m$ に当たる。平面波1つの位相速度は $\omega/k=\hbar k/(2m)$ なので，$v_{\mathrm{g}}$ はその約2倍である。

\shoumonkai{3} $\theta$ が実数なら $\abs{\ee^{-\ii\theta}}=1$ なので，$\abs{\psi_{n}}^{2}=\phi_{n}^{2}=\sin^{2}(n\pi x/L)$ で時間によらない。$\psi=\psi_{1}+\psi_{2}$ では，$\phi_{n}$ が実数であることを使って
\begin{align*}
  \abs{\psi}^{2}&=\phi_{1}^{2}+\phi_{2}^{2}+\phi_{1}\phi_{2}\left(\ee^{-\ii(E_{2}-E_{1})t/\hbar}+\ee^{\ii(E_{2}-E_{1})t/\hbar}\right)\\
  &=\sin^{2}\frac{\pi x}{L}+\sin^{2}\frac{2\pi x}{L}+2\sin\frac{\pi x}{L}\sin\frac{2\pi x}{L}\cos\frac{(E_{2}-E_{1})t}{\hbar}.
\end{align*}
$\abs{\psi}^{2}$ は角振動数 $\Omega=(E_{2}-E_{1})/\hbar$ で時間変化する。振動数は $\Omega/(2\pi)=(E_{2}-E_{1})/h$ で，$n=2$ から $n=1$ への遷移で放出される光子の振動数 $\nu$（$h\nu=E_{2}-E_{1}$）と一致する。
\kotae{\shoumonkai{1} $\omega=\frac{\hbar k^{2}}{2m}$（$E=p^{2}/(2m)$ を表す） \shoumonkai{2} $\abs{\psi}^{2}=2+2\cos[(k_{2}-k_{1})x-(\omega_{2}-\omega_{1})t]$，$v_{\mathrm{g}}=\frac{\hbar(k_{1}+k_{2})}{2m}$（2つの粒子の速さの平均），位相速度の約2倍 \shoumonkai{3} $\abs{\psi_{n}}^{2}=\sin^{2}(n\pi x/L)$；$\abs{\psi}^{2}=\sin^{2}(\pi x/L)+\sin^{2}(2\pi x/L)+2\sin(\pi x/L)\sin(2\pi x/L)\cos[(E_{2}-E_{1})t/\hbar]$，角振動数 $(E_{2}-E_{1})/\hbar$，振動数は光子の振動数 $(E_{2}-E_{1})/h$ に等しい}
\kenzan 次元：$\hbar k^{2}/m$ は $(\unit{\kilogram\metre\squared\per\second})(\unit{\per\metre\squared})/\unit{\kilogram}=\unit{\per\second}$。$k_{2}\to k_{1}$ で $v_{\mathrm{g}}\to\hbar k_{1}/m$ となり，粒子の速さに一致する。$\abs{\psi}^{2}$ は0から4の間の値をとる。\shoumonkai{3} で時間平均をとると交差項が消え，$\phi_{1}^{2}+\phi_{2}^{2}$ が残る。$x=L/2$ では $\phi_{2}=0$ なので，交差項はそこで振動しない。
\tsumazuki 定常状態1つの位相因子 $\ee^{-\ii E_{n}t/\hbar}$ は $\abs{\psi_{n}}^{2}$ に現れない。時間変化が見えるのは，エネルギーの異なる状態を重ね合わせたときだけである。
\end{kaisetsu}

\end{document}
